\begin{abstract}
Every complex quadratic pencil is recovered from its unmarked,
ungraded finite multiplication-failure algebra at a sharp order.
We relate this inverse to divisor-incidence geometry.  The entire
quadratic coefficient section of every multivariate reciprocal
normalization fibre is an apolar algebra of a power of the symmetric
determinant; its character, length, socle and nilpotence filtration
are explicit.  In one such section, the power-zero ideals along an
embedded pencil line are exactly powers of all spectral Fitting
ideals.  Their local lengths recover the Segre symbol and express
the native fixed-spectrum specialization order.  We also retain the
finite-flat binary factorization theorem, the squarefree conductor
and branch stratification, the complete incidence atlas for binary
collisions, and the full-coordinate-change common boundary theorem.
The latter boundary and the native pencil-line construction are
different objects.
\end{abstract}
\maketitle

\section{Introduction}
\label{sec:introduction-v153}
A finite multiplication-failure scheme may remember a multiplication
law even after its grading, coordinates and tensor factors have been
forgotten.  For quadratic pencils this recovery is exact at one sharp
infinitesimal order.  The first relation of the finite algebra also
has a common-divisor incidence.  We study how the nonreduced fibres
of that incidence record algebraic information and how a single
reciprocal fibre can carry the spectral schemes of a recovered pencil.

Let $\dim V=n\geq3$, $N=\binom{n+1}{2}$,
$R\in\Gr(2,\Sym^2V)$, $p=N-2$, and $D=n+2p=n^2+2n-4$.
For the universal matrix $T$ and quotient $\gamma_R$ by $R$, set
\[
 \mathfrak A_R^{[D]}=\C[t_{ij}]/
 \bigl((\det T)I_p(\gamma_R\Sym^2T)+\mathfrak m^{D+1}\bigr).
\]
Theorems~\ref{thm:artin-local-inverse-v146} and
\ref{thm:sharp-finite-pencil} recover the congruence class of every
pencil $R$ from the unmarked, ungraded algebra
$\mathfrak A_R^{[D]}$.  Every smaller truncation is independent of
$R$.  No regularity hypothesis on the pencil is imposed.  The
relative statements recover actual moving subbundles and source
bundles.  The coefficient line of the inverse also determines the
polarized moduli completion of
Theorem~\ref{thm:polarized-completion-v154}.

\subsection{A whole section of every reciprocal fibre}
The normalization of a common-divisor image parametrizes a marked
factor and the residual relation space.  At a pure-power common
factor its scheme fibre is the reciprocal algebra $F_{g,k}(W)$.
For $g\geq k\geq2$, write $d=\dim W$ and $h=\lfloor g/2\rfloor$.
Theorem~\ref{thm:quadratic-section-v155} identifies its quadratic
coefficient section, as an algebra rather than only as a tangent
space, with
\[
 B_{d,h}=\Sym((\Sym^2W)^*)/(\operatorname{coeff}Q^{h+1}).
\]
This is the apolar algebra of $(\det X)^h$ for a generic symmetric
matrix $X$.  Its degree-$p$ piece is
$\bigoplus_{\lambda\subseteq(h^d),\,|\lambda|=p}
\mathbb S_{2\lambda}(W^*)$.  The resulting product formula for its
length gives a uniform lower bound on the entire reciprocal fibre;
its socle gives a nonzero power of the maximal ideal of that fibre.
In the actual ternary-pencil common boundary this yields
\[
             \length F_{6,2}(\C^8)\geq922268360,
                     \qquad\mathfrak m^{24}\ne0.
\]
The number is the exact length of the section, not an asserted
length of the full fibre.  Gorenstein duality and the Lefschetz
property are likewise asserted for the section only.

\subsection{Spectral schemes from multiplication in one fibre}
Take the normalization fibre $F_{2h,2}(V^*)$ above the actual
incidence point $a^{2h+2}(\C a\oplus V^*)$ and its quadratic
section $B_{n,h}$.  Since $(B_{n,h})_1=\Sym^2V$, the pencil
embeds its line $\Lambda=\Pj(R)$ in this one algebra's degree-one
space.  Let $\mathcal J_p$ be the coefficient ideal of the
universal $p$th power on $\Lambda$, and let $\mathcal T_R$ be the
cokernel of its universal symmetric matrix.
Theorem~\ref{thm:power-Fitting-v155} proves the ideal identity
\[
              \mathcal J_{h(n-j)}=\Fitt_j(\mathcal T_R)^h
                           \quad(0\leq j\leq n).
\]
The equality includes singular pencils and nonreduced spectral
schemes.  Its proof computes the valuation of every power ideal
from the symmetric elementary exponents over a discrete valuation
ring; it does not pass from equality of supports to equality of
schemes.  For a regular pencil the successive local lengths, divided
by $h$, recover every block in the Segre symbol, with its common
projective spectral support.  The fixed-spectrum full-orthogonal
closure order is therefore expressed by inclusions of these
power-zero divisors.  The closure theorem itself is the inherited
Theorem~\ref{thm:fixed-spectral-classification}; the new point is its
realization by multiplication in one reciprocal section.

\subsection{The incidence geometry supporting the construction}
For binary forms the normalization is finite flat of degree
$\binom{s}{g}$ over every exact-common-divisor-degree-$s$ stratum.
Theorem~\ref{thm:binary-fibres-v153} determines every scheme fibre
as products of reciprocal complete intersections.
On the squarefree open $\mathcal U$,
Theorem~\ref{thm:squarefree-v153} gives
\[
 \mathfrak c_{Y_g}|_{\mathcal U}
                   =\II_{Y_{g+1}/Y_g}|_{\mathcal U},
\]
together with every multiple branch intersection, seminormality,
depth and multiplicity.  Theorem~\ref{thm:block-betti-v154}
computes the whole graded Betti table of its block model.  The
all-collision theorem remains an exact completed incidence atlas,
not an analytic classification of all collision singularities.
The double-root singular Fitting ideal, multivariate minimal relation
module, and complete five-variable resolution are retained as
Theorems~\ref{thm:fitting-v153},
\ref{thm:reciprocal-structure-v153} and
Proposition~\ref{prop:F22-v154}.

\subsection{Groups, classical inputs, and organization}
The power-zero construction uses the recovered pencil line and its
native congruence action.  It must not be confused with the common
extreme limit in all full $\operatorname{GL}(E)$ first-relation
orbit closures.  That limit is reached by two successive
specializations and does not remember the starting pencil.  Its
large reciprocal fibre supplies the length application above; the
embedded pencil line supplies the separate spectral application.
The first relation's original gcd remains $(\det T)^{n-1}$.
No equality with the auxiliary spectral ideals is asserted.

The determinant-power apolar representation and strong Lefschetz
theorems are classical results of \cite{NagaokaWachi2024}; for the
first power see also \cite{ShafieiSymmetric}.  We use their exact
ideal generators to identify the entire coefficient section and then
prove the power-ideal identity on pencil lines.  Binary reciprocal,
block-monomial and orbit-order inputs are credited to
\cite{Grinberg2019,MillerSturmfels2005,FevolaMandelshtamSturmfels}
and to the orthogonal-orbit comparison in the proof of
Theorem~\ref{thm:fixed-spectral-classification}.
The theorem and proof text of \cite{Ballico93} has not been obtained;
no theorem-level anticipation or nonanticipation assertion is made.

The arguments are also organized as two companion papers.  The first
contains the sharp inverse, its relative forms and the native moduli
and specialization theory.  The second contains the divisor-image
geometry, reciprocal algebra and power-zero spectral construction.
Their numerical references are explicitly distinguished by I and II.
This complete text retains both bodies of proofs, including the
recognition, curve, covering and rigidification arguments.

\part*{Part I. Reconstruction and intrinsic moduli}
